%%=============================================================================
%% Samenvatting
%%=============================================================================

% TODO: De "abstract" of samenvatting is een kernachtige (~ 1 blz. voor een
% thesis) synthese van het document.
%
% Een goede abstract biedt een kernachtig antwoord op volgende vragen:
%
% 1. Waarover gaat de bachelorproef?
% 2. Waarom heb je er over geschreven?
% 3. Hoe heb je het onderzoek uitgevoerd?
% 4. Wat waren de resultaten? Wat blijkt uit je onderzoek?
% 5. Wat betekenen je resultaten? Wat is de relevantie voor het werkveld?
%
% Daarom bestaat een abstract uit volgende componenten:
%
% - inleiding + kaderen thema
% - probleemstelling
% - (centrale) onderzoeksvraag
% - onderzoeksdoelstelling
% - methodologie
% - resultaten (beperk tot de belangrijkste, relevant voor de onderzoeksvraag)
% - conclusies, aanbevelingen, beperkingen
%
% LET OP! Een samenvatting is GEEN voorwoord!

%%---------- Nederlandse samenvatting -----------------------------------------
%
% TODO: Als je je bachelorproef in het Engels schrijft, moet je eerst een
% Nederlandse samenvatting invoegen. Haal daarvoor onderstaande code uit
% commentaar.
% Wie zijn bachelorproef in het Nederlands schrijft, kan dit negeren, de inhoud
% wordt niet in het document ingevoegd.

\IfLanguageName{english}{%
\selectlanguage{dutch}
\chapter*{Samenvatting}
\lipsum[1-4]
\selectlanguage{english}
}{}

%%---------- Samenvatting -----------------------------------------------------
% De samenvatting in de hoofdtaal van het document

\chapter*{\IfLanguageName{dutch}{Samenvatting}{Abstract}}

De beveiliging van beheertoegang tot servers vormt een kritieke uitdaging binnen moderne IT-omgevingen. Systeem- en netwerkbeheerders maken dagelijks gebruik van protocollen zoals Secure Shell (SSH) en Remote Desktop Protocol (RDP) om systemen met verhoogde rechten te beheren. Hoewel deze protocollen cryptografisch beveiligd zijn, bieden zij onvoldoende controle op het gebruik van privileged accounts, wat kan leiden tot misbruik, privilege-escalatie en laterale beweging binnen het netwerk. Gecompromitteerde beheerdersaccounts vormen bovendien een frequente aanvalsvector binnen hedendaagse cyberincidenten.

Deze bachelorproef onderzoekt hoe een Privileged Access Management (PAM)-tool kan bijdragen aan de beveiliging van servertoegang via SSH en RDP binnen een bedrijfsnetwerk. Het doel van het onderzoek is na te gaan in welke mate PAM-oplossingen beheertoegang beter kunnen controleren, tijdelijk kunnen toekennen en volledig kunnen loggen, zodat beveiligingsrisico’s worden beperkt en compliancevereisten worden ondersteund. Daarnaast wordt onderzocht hoe enterprise-organisaties PAM-capabilities prioriteren op basis van governance, operationele efficiëntie, gebruikservaring en beveiligingsrisico’s.

Het onderzoek werd uitgevoerd in twee fasen. Eerst werd een literatuurstudie en marktanalyse uitgevoerd op basis van toonaangevende analistenrapporten, waaronder het Gartner Magic Quadrant, Gartner Voice of the Customer, de KuppingerCole Leadership Compass en The Forrester Wave. Daarnaast werden twee casestudies uitgewerkt binnen Ahold Delhaize. De eerste casestudie behandelt de strategische heroriëntatie van de PAM-omgeving binnen de organisatie en onderzoekt hoe capabilities zoals governance, Just-in-Time toegang, credential vaulting, automatisatie en compliance prioritair werden opgenomen binnen de PAM-roadmap. De tweede casestudie focust op de operationele implementatie van CyberArk binnen een storageomgeving en vertaalt deze strategische keuzes naar concrete technische vereisten en dagelijkse beheerprocessen.

Op basis van de literatuurstudie, marktanalyse en casestudies werd een vergelijkende studie opgesteld van vier representatieve PAM-oplossingen. Vervolgens werd een proof of concept uitgevoerd op fysieke machines waarop een gesimuleerd bedrijfsnetwerk werd opgezet. In deze testomgeving werden twee PAM-oplossingen geïmplementeerd en geëvalueerd op basis van vooraf vastgelegde criteria, waaronder Just-in-Time toegang, sessielogging, credential vaulting en het blokkeren van directe verbindingen.

De resultaten tonen aan dat alle onderzochte PAM-oplossingen voldoen aan de basisvereisten voor het beveiligen van privileged access. Tegelijkertijd wordt een duidelijk onderscheid zichtbaar in maturiteit en functionaliteit. Twee van de vier oplossingen bieden een aanzienlijk hoger niveau van automatisatie, granulariteit en controle, waardoor ze beter geschikt zijn voor complexe en streng gereguleerde omgevingen. De andere twee oplossingen focussen eerder op eenvoud en snelle implementatie, en voldoen aan de basisbehoeften, maar bieden minder diepgang en flexibiliteit.

Daarnaast blijkt dat de effectiviteit van een PAM-oplossing niet enkel afhangt van de beschikbare functionaliteiten, maar in sterke mate bepaald wordt door de configuratie, governance en integratie binnen de bestaande IT-omgeving. De casestudies tonen bovendien aan dat moderne enterprise-organisaties PAM niet uitsluitend beschouwen als een oplossing voor password management, maar als een strategisch platform voor gecontroleerde toegang, compliance, auditing en operationele standaardisatie binnen hybride IT-omgevingen.

